\documentclass[11pt, a4paper]{article}

% Paquetes básicos
\usepackage[utf8]{inputenc}
\usepackage[T1]{fontenc}
\usepackage{lmodern}
\usepackage[spanish]{babel}
\usepackage[margin=0.8cm, landscape]{geometry}
\usepackage{multicol}
\usepackage{tcolorbox}
\usepackage{tabularx}
\usepackage{booktabs}
\usepackage{amsmath}

\title{\textbf{Formulario del Capítulo 1: Abstracciones y Tecnología de los Computadores}}
\author{Ronald Montilla}
\date{\today}

\begin{document}
{\centering \Large \textbf{Formulario Capítulo 1 - Estructura y Diseño de Computadores} \\[0.2cm]}
{\centering \normalsize \textbf{Basado en el libro de Patterson y Hennessy} \\ \vspace{0.4cm}}

    \begin{multicols*}{2}

        \begin{tcolorbox}[
        colback=white, 
        colframe=blue!40!black, 
        title= {\centering \textbf{Prestaciones y Tiempo de Ejecución} },
        fonttitle=\bfseries,
        size=small
        ]
        \begin{itemize}
            \item \textbf{Prestaciones:}
            $$\text{Prestaciones}_X = \frac{1}{\text{Tiempo de ejecución}_X}$$
            \item \textbf{Rendimiento relativo (X es $n$ veces más rápida que Y):}
            $$n = \frac{\text{Prestaciones}_X}{\text{Prestaciones}_Y} = \frac{\text{Tiempo de ejecución}_Y}{\text{Tiempo de ejecución}_X}$$
        \end{itemize}
        \end{tcolorbox}

        \begin{tcolorbox}[
        colback=white, 
        colframe=blue!40!black, 
        title= {\centering \textbf{Tiempo de CPU} },
        fonttitle=\bfseries,
        size=small
        ]
        \begin{itemize}
            \item \textbf{Basado en tiempo de ciclo:}
            $$\text{Tiempo de CPU} = \text{Ciclos de reloj de CPU} \times \text{Tiempo de ciclo de reloj}$$
            \item \textbf{Basado en frecuencia:}
            $$\text{Tiempo de CPU} = \frac{\text{Ciclos de reloj de CPU}}{\text{Frecuencia de reloj}}$$
        \end{itemize}
        \end{tcolorbox}

        \begin{tcolorbox}[
        colback=white, 
        colframe=blue!40!black, 
        title= {\centering \textbf{Ciclos de Reloj e Instrucciones (CPI)} },
        fonttitle=\bfseries,
        size=small
        ]
        \begin{itemize}
            \item \textbf{Ciclos de CPU totales:}
            $$\text{Ciclos de CPU} = \text{Nº de instrucciones} \times \text{CPI medio}$$
            \item \textbf{Ciclos basados en clases de instrucciones:}
            $$\text{Ciclos de CPU} = \sum_{i=1}^{n} (\text{CPI}_i \times C_i)$$
            \textit{Donde $C_i$ es el número de instrucciones de la clase $i$.}
            \item \textbf{CPI (Ciclos de reloj Por Instrucción):}
            $$\text{CPI} = \frac{\text{Ciclos de reloj de la CPU}}{\text{Número de instrucciones}}$$
        \end{itemize}
        \end{tcolorbox}

        \begin{tcolorbox}[
        colback=white, 
        colframe=blue!40!black, 
        title= {\centering \textbf{Ecuación Clásica de las Prestaciones} },
        fonttitle=\bfseries,
        size=small
        ]
        \begin{itemize}
            \item \textbf{Con tiempo de ciclo:}
            $$\text{Tiempo} = \text{Nº Instrucciones} \times \text{CPI} \times \text{Tiempo de ciclo}$$
            \item \textbf{Con frecuencia de reloj:}
            $$\text{Tiempo} = \frac{\text{Nº Instrucciones} \times \text{CPI}}{\text{Frecuencia de reloj}}$$
            \item \textbf{Desglose por componentes de las prestaciones:}
            $$\frac{\text{Segundos}}{\text{Programa}} = \frac{\text{Instrucciones}}{\text{Programa}} \times \frac{\text{Ciclos de reloj}}{\text{Instrucción}} \times \frac{\text{Segundos}}{\text{Ciclo de reloj}}$$
        \end{itemize}
        \end{tcolorbox}

        \vfill\null\columnbreak

        \begin{tcolorbox}[
        colback=white, 
        colframe=blue!40!black, 
        title= {\centering \textbf{Consumo de Potencia} },
        fonttitle=\bfseries,
        size=small
        ]
        \begin{itemize}
            \item \textbf{Potencia dinámica (tecnología CMOS):}
            $$\text{Potencia} = \text{Carga capacitiva} \times \text{Voltaje}^2 \times \text{Frecuencia de conmutación}$$
            \item \textbf{Potencia estática (por fugas de corriente):}
            $$P_{est} = V \times I_{\text{fuga}}$$
            \item \textbf{Métrica global SPECpower (ssj\_ops global por vatio):}
            $$\text{Rendimiento/Potencia} = \frac{\sum_{i=0}^{10} \text{ssj\_ops}_i}{\sum_{i=0}^{10} \text{potencia}_i}$$
        \end{itemize}
        \end{tcolorbox}

        \begin{tcolorbox}[
        colback=white, 
        colframe=blue!40!black, 
        title= {\centering \textbf{Ley de Amdahl} },
        fonttitle=\bfseries,
        size=small
        ]
        \begin{itemize}
            \item \textbf{Tiempo de ejecución después de las mejoras:}
            $$T_{\text{nuevo}} = \frac{T_{\text{afectado por mejora}}}{\text{Cantidad de mejora}} + T_{\text{no afectado}}$$
            \textit{Nota: Sirve para calcular el impacto de mejorar sólo una fracción del sistema.}
        \end{itemize}
        \end{tcolorbox}

        \begin{tcolorbox}[
        colback=white, 
        colframe=blue!40!black, 
        title= {\centering \textbf{Coste de Fabricación de Circuitos Integrados} },
        fonttitle=\bfseries,
        size=small
        ]
        \begin{itemize}
            \item \textbf{Coste por dado (chip):}
            $$\text{Coste} = \frac{\text{Coste por oblea}}{\text{Dados por oblea} \times \text{Factor de producción}}$$
            \item \textbf{Dados por oblea (aproximación):}
            $$\text{Dados por oblea} \approx \frac{\text{Área de la oblea}}{\text{Área del dado}}$$
            \item \textbf{Factor de producción (Yield):}
            $$\text{Factor} = \frac{1}{\left(1 + \left(\text{Defectos por área} \times \frac{\text{Área del dado}}{2}\right)\right)^2}$$
        \end{itemize}
        \end{tcolorbox}
        
    \end{multicols*}

\end{document}